\section{Empirical Evaluation}
\label{sec:diagnosis}

To demonstrate the framework's application, we evaluate systems across three categories: government infrastructure promoted as ``DPI,'' platform infrastructure claiming ``openness,'' and infrastructure that satisfies all five tests.

\subsection{Government Infrastructure: The Trap of Partial Compliance}

Government-operated systems currently designated as Digital Public Infrastructure frequently exhibit a dangerous failure mode: they achieve partial compliance on technical tests (EXIT, CODE, AUDIT) while failing structurally on governance tests (GOVERN, FORK). By capturing the monopoly on execution, they render the feedback loop inoperable.

\begin{table}[htbp]
    \centering
    \caption{\textbf{Evaluation of Government Infrastructure.} Centralized execution leads to consistent structural failure despite ``open standards'' rhetoric.}
    \label{tab:govt_infra}

    % --- Subtable A: Aadhaar ---
    \begin{subtable}{1\textwidth}
        \centering
        \caption{Aadhaar (India) -- Mandatory ID}
        \small
        \begin{tabular}{@{}llp{8.5cm}@{}}
            \toprule
            \textbf{Test} & \textbf{Result} & \textbf{Evidence} \\
            \midrule
            EXIT & \textcolor{red}{FAIL} & Mandatory for essential survival services (PDS, Banking) \citep{khera2019}. \\
            CODE & \textcolor{red}{FAIL} & Proprietary matcher. Execution logic is opaque. \\
            AUDIT & \textcolor{red}{FAIL} & Independent failure-rate testing prohibited by law. \\
            GOVERN & \textcolor{red}{FAIL} & UIDAI (Operator) acts as Regulator. No citizen veto. \\
            FORK & \textcolor{red}{FAIL} & Total centralized capture; cannot be reproduced. \\
            \bottomrule
        \end{tabular}
        \vspace{0.3cm}
    \end{subtable}

    % --- Subtable B: UPI ---
    \begin{subtable}{1\textwidth}
        \centering
        \caption{UPI (India) -- Payment Rail}
        \small
        \begin{tabular}{@{}llp{8.5cm}@{}}
            \toprule
            EXIT & \textcolor{orange}{PARTIAL} & Cash exists but network effects coerce adoption. \\
            CODE & \textcolor{orange}{PARTIAL} & Open protocol specs; Closed switching logic. \\
            AUDIT & \textcolor{orange}{PARTIAL} & Regulator audits allowed; Public verification blocked. \\
            GOVERN & \textcolor{red}{FAIL} & Bank-controlled governance (NPCI). No user seat. \\
            FORK & \textcolor{red}{FAIL} & Hub-and-spoke monopoly cannot be forked. \\
            \bottomrule
        \end{tabular}
        \vspace{0.3cm}
    \end{subtable}

    % --- Subtable C: Pix ---
    \begin{subtable}{1\textwidth}
        \centering
        \caption{Pix (Brazil) -- Central Bank Payment}
        \small
        \begin{tabular}{@{}llp{8.5cm}@{}}
            \toprule
            EXIT & \textcolor{orange}{PARTIAL} & Other payment methods exist; increasing state reliance. \\
            CODE & \textcolor{orange}{PARTIAL} & Protocol public; Central switch proprietary. \\
            AUDIT & \textcolor{orange}{PARTIAL} & Volume data public; Granular error data private. \\
            GOVERN & \textcolor{red}{FAIL} & Unilateral Central Bank governance. \\
            FORK & \textcolor{red}{FAIL} & Cannot be reproduced independently of BCB. \\
            \bottomrule
        \end{tabular}
    \end{subtable}
\end{table}

\textbf{Pattern}: All three systems fail GOVERN and FORK. Partial compliance on technical tests does not satisfy the binary threshold.

\subsection{Platform Infrastructure: The Illusion of Openness}

In the private sector, systems often leverage "open source" branding while retaining centralized control. This decoupling of artifact transparency from governance reality creates a distinct failure pattern: the artifacts are public; the decisions are not.

\begin{table}[htbp]
    \centering
    \caption{\textbf{Evaluation of Private/Platform Infrastructure.} Technical openness (Source/Weights) does not guarantee structural accountability.}
    \label{tab:corp_infra}

    % --- Subtable A: Open Weights AI ---
    \begin{subtable}{1\textwidth}
        \centering
        \caption{Open-Weights AI (Llama, etc.)}
        \small
        \begin{tabular}{@{}llp{8.5cm}@{}}
            \toprule
            \textbf{Test} & \textbf{Result} & \textbf{Evidence} \\
            \midrule
            EXIT & \textcolor{green!50!black}{PASS} & No vendor lock-in; can run locally. \\
            CODE & \textcolor{orange}{PARTIAL} & Weights public; Training Data/RLHF opaque. \\
            AUDIT & \textcolor{orange}{PARTIAL} & Black-box testing possible; Lineage unverifiable. \\
            GOVERN & \textcolor{red}{FAIL} & Safety/Utility trade-offs set unilaterally by firm. \\
            FORK & \textcolor{orange}{PARTIAL} & Fine-tune possible; Retraining blocked by resource cost. \\
            \bottomrule
        \end{tabular}
        \vspace{0.3cm}
    \end{subtable}

    % --- Subtable B: Android ---
    \begin{subtable}{1\textwidth}
        \centering
        \caption{Android (AOSP + GMS)}
        \small
        \begin{tabular}{@{}llp{8.5cm}@{}}
            \toprule
            EXIT & \textcolor{orange}{PARTIAL} & Ecosystem lock-in high despite device portability. \\
            CODE & \textcolor{orange}{PARTIAL} & OS open source; Identity/Push services proprietary. \\
            AUDIT & \textcolor{orange}{PARTIAL} & OS auditable; Google Services opaque. \\
            GOVERN & \textcolor{red}{FAIL} & Roadmap dictated by Google. \\
            FORK & \textcolor{orange}{PARTIAL} & LineageOS possible; App ecosystem breaks without GMS. \\
            \bottomrule
        \end{tabular}
        \vspace{0.3cm}
    \end{subtable}

    % --- Subtable C: Signal ---
    \begin{subtable}{1\textwidth}
        \centering
        \caption{Signal Protocol}
        \small
        \begin{tabular}{@{}llp{8.5cm}@{}}
            \toprule
            EXIT & \textcolor{green!50!black}{PASS} & No lock-in. \\
            CODE & \textcolor{green!50!black}{PASS} & Fully open source (Client/Server). \\
            AUDIT & \textcolor{green!50!black}{PASS} & Cryptographically verifiable. \\
            GOVERN & \textcolor{red}{FAIL} & Foundation acts unilaterally. \\
            FORK & \textcolor{red}{FAIL} & Code forkable; Network locked (no federation). \\
            \bottomrule
        \end{tabular}
    \end{subtable}
\end{table}

\textbf{Pattern}: Open artifacts with closed governance. ``Open weights'' $\neq$ ``open.'' ``Open source'' $\neq$ ``accountable.''

\subsection{Corrigible Infrastructure}

For contrast, Table~\ref{tab:positive} summarizes seventeen systems that pass all five tests.

\begin{table}[htbp]
\centering
\small
\renewcommand{\arraystretch}{1.3}
\begin{tabular}{@{}l>{\raggedright}p{3.2cm}>{\centering\arraybackslash}p{1.2cm}p{5.8cm}@{}}
\toprule
\rowcolor{gray!15}
\textbf{System} & \textbf{Steward} & \textbf{Score} & \textbf{Why Corrigible} \\
\midrule
\rowcolor{green!8}
Linux Kernel & Linux Foundation & \cellcolor{green!25}\textbf{5/5} & Fully reproducible. Full triad (legal, technical, economic) exists. Independent distributions prove forkability. \\
Let's Encrypt & ISRG & \cellcolor{green!25}\textbf{5/5} & Open source CA, IETF standards, nonprofit 501(c)(3) \\
\rowcolor{green!8}
Wikipedia & Wikimedia Foundation & \cellcolor{green!25}\textbf{5/5} & CC BY-SA content, public edit history, community RfC governance \\
Matrix Protocol & Matrix.org Foundation & \cellcolor{green!25}\textbf{5/5} & Apache 2.0, federated by design, self-hostable \\
\rowcolor{green!8}
Bluesky (AT Protocol) & Bluesky PBC & \cellcolor{green!25}\textbf{5/5} & MIT/Apache, portable DID identity, self-hostable PDS \\
PostgreSQL & PGDG & \cellcolor{green!25}\textbf{5/5} & BSD-like license, public mailing lists, major forks exist \\
\rowcolor{green!8}
IPFS & Protocol Labs & \cellcolor{green!25}\textbf{5/5} & MIT/Apache, public IPIP process, no token governance \\
Bitcoin & Decentralized & \cellcolor{green!25}\textbf{5/5} & MIT source, public BIP process, proven forks (BCH, BSV) \\
\rowcolor{green!8}
Kubernetes & CNCF & \cellcolor{green!25}\textbf{5/5} & Apache 2.0, KEP process, distributions (OpenShift, k3s) \\
Firefox & Mozilla Foundation & \cellcolor{green!25}\textbf{5/5} & MPL 2.0, nonprofit manifesto, forks (LibreWolf, Tor Browser) \\
\rowcolor{green!8}
Apache HTTP & Apache Foundation & \cellcolor{green!25}\textbf{5/5} & Apache 2.0 license, ASF governance, powers 30\% of web \\
Apache Kafka & Apache Foundation & \cellcolor{green!25}\textbf{5/5} & Apache 2.0, KIP process public, Fortune 500 standard \\
\rowcolor{green!8}
OpenSearch & Linux Foundation & \cellcolor{green!25}\textbf{5/5} & Forked from Elasticsearch; Apache 2.0, multi-vendor governance \\
Valkey & Linux Foundation & \cellcolor{green!25}\textbf{5/5} & Forked from Redis; BSD license, community governance restored \\
\rowcolor{green!8}
Hyperledger & LF Decentralized Trust & \cellcolor{green!25}\textbf{5/5} & Apache 2.0 framework; institutional deployments require separate evaluation \\
LibreOffice & Document Foundation & \cellcolor{green!25}\textbf{5/5} & Forked from OpenOffice; LGPL/MPL, nonprofit governance \\
\rowcolor{green!8}
MariaDB & MariaDB Foundation & \cellcolor{green!25}\textbf{5/5} & Forked from MySQL; GPL, foundation governance \\
Eclipse IDE & Eclipse Foundation & \cellcolor{green!25}\textbf{5/5} & EPL 2.0, 386 member orgs, committer-led governance \\
\bottomrule
\end{tabular}
\caption{Infrastructure systems satisfying all five corrigibility tests. All systems demonstrate full EXIT, CODE, AUDIT, GOVERN, and FORK compliance.}
\label{tab:positive}
\end{table}

\textbf{Fork as Correction}: OpenSearch, Valkey, LibreOffice, and MariaDB demonstrate Ashby's Law in practice. Each emerged when governance variety collapsed: Oracle's Sun acquisition (2010) triggered LibreOffice and MariaDB; unilateral license changes triggered OpenSearch (2021) and Valkey (2024). The ecosystem restored requisite variety through forking.

The framework is governance-agnostic: corporate participation is not disqualifying. Kubernetes (CNCF), Hyperledger, and Let's Encrypt all have significant corporate involvement while passing all tests. The test is whether governance permits correction, not who participates.

\textbf{License vs. Governance}: Redis Ltd. added AGPLv3 in 2025, restoring open licensing. But license is orthogonal to governance. \textit{License} determines what you may do with code; \textit{governance} determines who decides what it becomes. Redis passes one; Valkey passes both.

\subsection{The Taxonomy of Open-Washing}
\label{subsec:openwashing}

Systems that fail these structural tests frequently employ vocabulary to obscure their closure:
\begin{enumerate}
    \item \textbf{Standard-Washing}: Adopting open technical standards (like ISO/W3C) as proxy for governance accountability.
    \item \textbf{Open-Washing}: Releasing peripheral SDKs while keeping core execution logic proprietary.
    \item \textbf{Multilateral-Washing}: Using international endorsements (G20, WB) to substitute for technical auditability.
    \item \textbf{Safety-Washing}: Justifying opacity by claiming inspection poses a security risk.
    \item \textbf{Consent-Washing}: Obtaining formal consent under conditions of economic coercion.
    \item \textbf{DID-Washing}: Deploying decentralized identifiers without decentralized governance (see Section~\ref{subsubsec:didwashing}).
\end{enumerate}

\subsubsection{DID-Washing: The Layer Confusion Problem}
\label{subsubsec:didwashing}

Decentralized Identifiers (DIDs) present a distinct challenge to corrigibility evaluation because they separate technical decentralization from governance decentralization across multiple architectural layers. The W3C DID Core 1.0 specification \citep{w3cdid2022} defines DIDs as ``globally unique identifiers'' that ``do not require a centralized registration authority.'' This technical claim is accurate. However, a system can implement DIDs at the identifier layer while retaining centralized control at every layer that determines actual utility.

Identity infrastructure comprises at least four functional layers, each with independent governance:

\begin{table}[htbp]
\centering
\small
\begin{tabular}{@{}llp{5cm}@{}}
\toprule
\textbf{Layer} & \textbf{Function} & \textbf{DID-Washing Failure Mode} \\
\midrule
Identifier & Generate and resolve DIDs & Often decentralized, but meaningless alone \\
Credential & Issue verifiable claims about the subject & Usually centralized: single issuer dominates \\
Trust & Determine which credentials are accepted & Controlled by policy, trusted lists, or law \\
Revocation & Invalidate credentials or identifiers & Often requires issuer cooperation \\
\bottomrule
\end{tabular}
\caption{DID-washing failure modes by layer}
\label{tab:didwashing}
\end{table}

A DID that resolves correctly but is not accepted by any relying party is functionally equivalent to no identity at all. The identifier is ``self-sovereign''; the identity utility is not.

\paragraph{Empirical Examples.} The EU Digital Identity Wallet (EUDI) under eIDAS 2.0 \citep{eidas2024} illustrates sophisticated DID-washing. The architecture employs selective disclosure credentials and cryptographic proofs superficially similar to DID/VC ecosystems. However:

\begin{itemize}[noitemsep,topsep=0pt]
    \item Credential issuance for Person Identification Data is an exclusive government function.
    \item Trust determination requires registration in Member State Trusted Lists.
    \item Wallet certification depends on government-approved Qualified Trust Service Providers (QTSPs).
    \item Trusted Execution Environment dependency hands cryptographic key governance to mobile OS vendors (Google, Apple).
\end{itemize}

The result: users control their container (the wallet) but not their capability (the credentials). A journalist critical of their government holds a ``decentralized'' identifier that can be rendered useless by a single administrative act: revocation from the trusted list.

\paragraph{Implications for Framework Application.} DID-washing exploits a gap in naive corrigibility evaluation. A system may appear to pass CODE (the DID spec is public) and FORK (anyone can implement the spec) while failing at the trust layer where actual power resides.

The framework must be applied layer by layer (see Section~\ref{subsec:layers}). For identity systems:

\begin{itemize}[noitemsep,topsep=0pt]
    \item CODE must be evaluated at each layer: identifier resolution, credential issuance logic, trust framework rules, and revocation mechanisms.
    \item GOVERN must ask: who controls the acceptance policy? If a single authority determines which credentials are valid, governance is centralized regardless of identifier architecture.
    \item FORK must ask: can a community reproduce the trust framework, not merely the identifier scheme? If relying parties are legally required to accept only government-issued credentials, forkability is illusory.
\end{itemize}

A system passes the DID-washing test only if decentralization extends through the full trust stack: issuance, verification, acceptance policy, and revocation governance. Technical decentralization at the identifier layer with governance centralization at the trust layer is structural theater.

\begin{claim}
Identifier decentralization is not system decentralization. The corrigibility of an identity system is determined by its weakest layer across any test dimension.
\end{claim}

\begin{figure}[htbp]
\centering
\begin{tikzpicture}[
    scale=0.85,
    pass/.style={fill=green!30},
    fail/.style={fill=red!30},
    partial/.style={fill=orange!30},
    header/.style={font=\scriptsize\bfseries},
    cell/.style={font=\scriptsize, minimum width=0.9cm, minimum height=0.5cm}
]
% Headers
\node[header] at (0,0) {};
\node[header, text=exitcolor] at (1.2,0) {EXIT};
\node[header, text=codecolor] at (2.2,0) {CODE};
\node[header, text=auditcolor] at (3.2,0) {AUDIT};
\node[header, text=governcolor] at (4.2,0) {GOV};
\node[header, text=forkcolor] at (5.2,0) {FORK};
\node[header] at (6.4,0) {Result};

% Aadhaar
\node[cell, anchor=east] at (-0.2,-0.6) {Aadhaar};
\node[cell, fail] at (1.2,-0.6) {\texttimes};
\node[cell, fail] at (2.2,-0.6) {\texttimes};
\node[cell, fail] at (3.2,-0.6) {\texttimes};
\node[cell, fail] at (4.2,-0.6) {\texttimes};
\node[cell, fail] at (5.2,-0.6) {\texttimes};
\node[cell, font=\scriptsize\bfseries, text=red] at (6.4,-0.6) {0/5};

% UPI
\node[cell, anchor=east] at (-0.2,-1.2) {UPI};
\node[cell, partial] at (1.2,-1.2) {\textasciitilde};
\node[cell, partial] at (2.2,-1.2) {\textasciitilde};
\node[cell, partial] at (3.2,-1.2) {\textasciitilde};
\node[cell, fail] at (4.2,-1.2) {\texttimes};
\node[cell, fail] at (5.2,-1.2) {\texttimes};
\node[cell, font=\scriptsize\bfseries, text=red] at (6.4,-1.2) {0/5};

% Pix
\node[cell, anchor=east] at (-0.2,-1.8) {Pix};
\node[cell, partial] at (1.2,-1.8) {\textasciitilde};
\node[cell, partial] at (2.2,-1.8) {\textasciitilde};
\node[cell, partial] at (3.2,-1.8) {\textasciitilde};
\node[cell, fail] at (4.2,-1.8) {\texttimes};
\node[cell, fail] at (5.2,-1.8) {\texttimes};
\node[cell, font=\scriptsize\bfseries, text=red] at (6.4,-1.8) {0/5};

% Open-weights AI
\node[cell, anchor=east] at (-0.2,-2.4) {Llama etc.};
\node[cell, pass] at (1.2,-2.4) {\checkmark};
\node[cell, partial] at (2.2,-2.4) {\textasciitilde};
\node[cell, partial] at (3.2,-2.4) {\textasciitilde};
\node[cell, fail] at (4.2,-2.4) {\texttimes};
\node[cell, partial] at (5.2,-2.4) {\textasciitilde};
\node[cell, font=\scriptsize\bfseries, text=red] at (6.4,-2.4) {1/5};

% Separator line
\draw[gray, thick] (-0.8,-2.9) -- (7,-2.9);

% Linux
\node[cell, anchor=east] at (-0.2,-3.4) {Linux};
\node[cell, pass] at (1.2,-3.4) {\checkmark};
\node[cell, pass] at (2.2,-3.4) {\checkmark};
\node[cell, pass] at (3.2,-3.4) {\checkmark};
\node[cell, pass] at (4.2,-3.4) {\checkmark};
\node[cell, pass] at (5.2,-3.4) {\checkmark};
\node[cell, font=\scriptsize\bfseries, text=green!50!black] at (6.4,-3.4) {5/5};

% Let's Encrypt
\node[cell, anchor=east] at (-0.2,-4.0) {Let's Enc.};
\node[cell, pass] at (1.2,-4.0) {\checkmark};
\node[cell, pass] at (2.2,-4.0) {\checkmark};
\node[cell, pass] at (3.2,-4.0) {\checkmark};
\node[cell, pass] at (4.2,-4.0) {\checkmark};
\node[cell, pass] at (5.2,-4.0) {\checkmark};
\node[cell, font=\scriptsize\bfseries, text=green!50!black] at (6.4,-4.0) {5/5};

% Kubernetes (CNCF)
\node[cell, anchor=east] at (-0.2,-4.6) {K8s};
\node[cell, pass] at (1.2,-4.6) {\checkmark};
\node[cell, pass] at (2.2,-4.6) {\checkmark};
\node[cell, pass] at (3.2,-4.6) {\checkmark};
\node[cell, pass] at (4.2,-4.6) {\checkmark};
\node[cell, pass] at (5.2,-4.6) {\checkmark};
\node[cell, font=\scriptsize\bfseries, text=green!50!black] at (6.4,-4.6) {5/5};

% Bitcoin (Decentralized)
\node[cell, anchor=east] at (-0.2,-5.2) {Bitcoin};
\node[cell, pass] at (1.2,-5.2) {\checkmark};
\node[cell, pass] at (2.2,-5.2) {\checkmark};
\node[cell, pass] at (3.2,-5.2) {\checkmark};
\node[cell, pass] at (4.2,-5.2) {\checkmark};
\node[cell, pass] at (5.2,-5.2) {\checkmark};
\node[cell, font=\scriptsize\bfseries, text=green!50!black] at (6.4,-5.2) {5/5};

% +15 more (reference)
\node[cell, anchor=east, font=\scriptsize\itshape, text=gray] at (-0.2,-5.8) {+15 more};
\node[cell, font=\scriptsize, text=gray] at (3.2,-5.8) {See Table~\ref{tab:positive}};

% Legend
\node[font=\tiny, gray] at (3.2,-6.4) {\checkmark = Pass \quad \textasciitilde\ = Partial \quad \texttimes\ = Fail};
\end{tikzpicture}
\caption{Comparative evaluation summary. Systems labeled ``DPI'' (above line) fail structural tests. Corrigible infrastructure (below line) passes all five. See Table~\ref{tab:positive} for complete list.}
\label{fig:comparison}
\end{figure}

\subsection{Implications}

The contrast is instructive. Aadhaar and UPI, promoted as model DPI, fail all tests. The seventeen systems in Table~\ref{tab:positive}, rarely described as ``DPI,'' pass all tests.

This suggests the ``public'' in Digital Public Infrastructure often refers to \textit{population-scale deployment} rather than \textit{public accountability}. The framework distinguishes these: scale is not governance.

All systems that pass share common properties:
\begin{itemize}
    \item Open source with permissive or copyleft licensing
    \item Nonprofit or community governance structures
    \item Protocol standardization through open processes (IETF, LKML)
    \item No monopoly on essential services
\end{itemize}

Scale does not require capture. Universality does not require opacity. The choice is architectural.